\section{Complete classification at the ancestor midpoint}
\label{sec:midpoint}

The decoder concerns columns. The original arithmetic question concerns
Gram matrices with the same first diagonal entry. The bridge between
them is a full matrix conjugacy, which we prove before invoking the word
theorem. This avoids any choice of a square root of a Gram matrix or of
a middle letter in an odd palindrome.

\begin{lemma}[The Gram bridge]\label{lem:gram-bridge}
Fix a directive $d$ and two palindromes $z_X,z_Y$. Write
$M_i=M(F_d(z_i))$ and $\rho_i=\rho(F_d(z_i))$.
The following conditions are equivalent:
\begin{align}
 M_X=M_Y=M,\qquad \rho_X+\rho_Y&=2r_0M;
 \label{eq:ancestor-midpoint-condition}\\
 \nu_d(z_Y)&=R_d\nu_d(z_X)R_d.
 \label{eq:full-reflection-conjugacy}
\end{align}
Under these conditions,
$\nu_d(z_Y)e=R_d\nu_d(z_X)e$.
\end{lemma}

\begin{proof}
For a symmetric determinant-one matrix
$\left(\begin{smallmatrix}M&\rho\\\rho&H\end{smallmatrix}\right)$,
conjugation by $J_0$ as a bilinear form gives
$$
 J_0^{\mathsf T}
 \begin{pmatrix}M&\rho\\\rho&H\end{pmatrix}J_0
 =\begin{pmatrix}
 M&2r_0M-\rho\\
 2r_0M-\rho&4r_0^2M-4r_0\rho+H
 \end{pmatrix}.
$$
The lower-right entry is uniquely determined by the other entries and
the determinant, because $M>0$. Hence
\eqref{eq:ancestor-midpoint-condition} is equivalent to
$\mc G(F_d(z_Y))=J_0^{\mathsf T}\mc G(F_d(z_X))J_0$.
Substitute \eqref{eq:image-gram}, multiply by $P^{-1}$ on both sides,
and use $R_d=PJ_0P^{-1}$ to obtain
$$
 G\nu_d(z_Y)=R_d^{\mathsf T}G\nu_d(z_X)R_d
 =GR_d\nu_d(z_X)R_d.
$$
Cancel the invertible matrix $G$. All steps are reversible, proving
the equivalence. Finally $R_de=e$ gives the column identity.
\end{proof}

\begin{theorem}[Palindromic pairs at the ancestor midpoint]
\label{thm:palindromic-classification}
Fix any finite directive $d$. Two binary palindromes $z_X,z_Y$
satisfy \eqref{eq:ancestor-midpoint-condition} if and only if, for
some word $w\in\{ab,ba\}^*$,
\begin{equation}\label{eq:classified-palindromes}
 z_X=\widetilde w w,
 \qquad z_Y=\widetilde{\widehat w}\,\widehat w.
\end{equation}
The word $w$ is uniquely determined by the ordered pair.
If it contains $k$ blocks, then
\begin{equation}\label{eq:classified-counts}
 c(F_d(z_X))=c(F_d(z_Y))
 =(2k+1)L_d\binom11.
\end{equation}
The case $k=0$ consists of two copies of the ancestor $g_d$.
Every case $k>0$ consists of two distinct, noncentral palindromes.
\end{theorem}

\begin{proof}
The Gram bridge and \cref{thm:word-decoder}, applied to the
whole words $z_X,z_Y$, show that
$z_X\in\{ab,ba\}^*$ and $z_Y=\widehat{z_X}$.
Because $z_X$ has even length and is a palindrome, each of its two
letter counts is even. A word of $m$ blocks $ab,ba$ has $m$ copies
of each letter. Thus $m=2k$ is even, and the midpoint of $z_X$
falls between blocks. Its second half is a unique word
$w\in\{ab,ba\}^*$ and $z_X=\widetilde w w$.
Block exchange commutes with reversal: reversing a block interchanges
$ab$ and $ba$, and reversing a sequence reverses the order of its
blocks. Consequently
$z_Y=\widetilde{\widehat w}\,\widehat w$.

Conversely, the block intertwining identities in
\eqref{eq:decoder-intertwiners} imply
$R_d\nu_d(x)R_d=\nu_d(\widehat x)$ for every
$x\in\{ab,ba\}^*$. Apply this with $x=\widetilde w w$ and then
use the Gram bridge. This proves the converse at the matrix level,
including the equality of labels.

Both preimages in \eqref{eq:classified-palindromes} have $2k$ copies
of each letter. Formula~\eqref{eq:completed-count-transport} proves
\eqref{eq:classified-counts}. If $k>0$, these completed counts have
greatest common divisor $2k+1>1$, whereas every central word has
primitive completed counts. The words are distinct because block
exchange changes their first letter, and $F_d$ is injective. The
empty case follows from $F_d(\eps)=g_d$.
\end{proof}

\begin{corollary}[Exclusion for central occurrences]
\label{cor:no-central-midpoint}
Let $d,t_X,t_Y$ be finite directives. If the central words
$\operatorname{Pal}(dt_X)$ and $\operatorname{Pal}(dt_Y)$ have the
same Gram label $M$ and their roots sum to $2M\rho_0/M_0$, then
$t_X=t_Y=\eps$. In particular, a pair of distinct central words
with equal labels cannot have its root midpoint at a common
standard ancestor's root parameter.
\end{corollary}

\begin{proof}
Use $z_i=\operatorname{Pal}(t_i)$ in the theorem. The images are
central, so only $k=0$ is possible. The palindromization recurrence
strictly increases length at every nonempty directive step; hence
$\operatorname{Pal}(t_i)=\eps$ implies $t_i=\eps$.
\end{proof}

\subsection{Exact examples and the role of primitivity}

At the empty directive, $G=I$, $(M_0,\rho_0)=(5,2)$ and
$$
 R_\eps=\frac15\begin{pmatrix}3&4\\4&-3\end{pmatrix}.
$$
Take $w=ab$. Then $z_X=baab$ and $z_Y=abba$, and direct evaluation gives
\begin{equation}\label{eq:small-gram-pair}
 \mc G(baab)=\begin{pmatrix}1130&467\\467&193\end{pmatrix},
 \qquad
 \mc G(abba)=\begin{pmatrix}1130&437\\437&169\end{pmatrix}.
\end{equation}
Both determinants are one and
$467+437=2\cdot1130\cdot(2/5)$.
This is a real equality of integral Gram labels. Its completed
counts are $(3,3)$, so the theorem explains exactly why it is
irrelevant to a collision between distinct central occurrences.

The same preimages give a pair for every directive, and the whole table
is one formula. It holds around every palindrome, standard or not.

\begin{table}[tb]
\centering
\caption{Images of the fixed palindromes $baab$ and $abba$. Here
$(M_0,\rho_0)$ belongs to the standard ancestor $\operatorname{Pal}(d)$,
$M$ is the common image label, and $\rho_-<\rho_+$ are the image
roots. Every row satisfies $(\rho_-+\rho_+)/(2M)=\rho_0/M_0$.
The completed counts have greatest common divisor three in every row.}
\label{tab:pairs}
\small
\setlength{\tabcolsep}{9pt}
\renewcommand{\arraystretch}{1.12}
\begin{tabular}{@{}crrrrr@{}}
\toprule
$d$ & $M_0$ & $\rho_0$ & $M$ & $\rho_-$ & $\rho_+$ \\
\midrule
$\eps$ & 5 & 2 & 1130 & 437 & 467 \\
$a$ & 13 & 5 & $19\,786$ & 7571 & 7649 \\
$b$ & 29 & 12 & $219\,530$ & $90\,753$ & $90\,927$ \\
$ab$ & 194 & 75 & $65\,712\,650$ & $25\,403\,793$ & $25\,404\,957$ \\
$ba$ & 433 & 179 & $730\,645\,066$ & $302\,043\,659$ & $302\,046\,257$ \\
\bottomrule
\end{tabular}
\end{table}

\begin{proposition}[The one-block pair in closed form]
\label{prop:one-block-closed-form}
Let $p$ be a palindrome with Gram label $M_p$ and Gram root $\rho_p$. The words
$p\,ba\,p\,ab\,p$ and $p\,ab\,p\,ba\,p$ have the common Gram label and the two Gram roots
\begin{equation}\label{eq:one-block-closed-form}
 M=M_p\bigl(9M_p^2+1\bigr),
 \qquad
 \rho_\pm=\bigl(9M_p^2+1\bigr)\rho_p\pm3M_p ,
\end{equation}
the sign being $+$ for $p\,ba\,p\,ab\,p$ and $-$ for $p\,ab\,p\,ba\,p$. In particular the two
roots differ by exactly $6M_p$, and their mean is $M\rho_p/M_p$. For $p=\operatorname{Pal}(d)$
the two words are $F_d(baab)$ and $F_d(abba)$, and $(M_p,\rho_p)=(M_0,\rho_0)$.
\end{proposition}

\begin{proof}
Put $W=\mu(p)$ and $\mc G_p=PWP=\left(\begin{smallmatrix}m&r\\r&h\end{smallmatrix}\right)$, a
symmetric matrix with $mh-r^2=1$, where $m=M_p$ and $r=\rho_p$. Then
$\mu(p\,ba\,p\,ab\,p)=W\,BA\,W\,AB\,W$, and inserting $P^{-1}P$ between the factors gives
$$
 \mc G(p\,ba\,p\,ab\,p)=\mc G_p\,K\,\mc G_p\,K^{\mathsf T}\mc G_p,
 \qquad
 K=P^{-1}BAP^{-1}=\begin{pmatrix}3&-1\\1&0\end{pmatrix},
$$
where $K^{\mathsf T}=P^{-1}ABP^{-1}$ because $A$, $B$ and $P$ are symmetric. A direct multiplication gives
$$
 K\mc G_pK^{\mathsf T}=\begin{pmatrix}9m-6r+h&3m-r\\3m-r&m\end{pmatrix},
 \qquad
 \mc G_pK\mc G_pK^{\mathsf T}=\begin{pmatrix}9m^2-3mr+1&3m^2\\9mr-3r^2+3&3mr+1\end{pmatrix},
$$
where $mh=r^2+1$ was used in the second product. Multiplying by the first column
$(m,r)^{\mathsf T}$ of $\mc G_p$, the first column of the Gram matrix is
$$
 \binom{(9m^2-3mr+1)m+3m^2r}{(9mr-3r^2+3)m+(3mr+1)r}
 =\binom{m(9m^2+1)}{(9m^2+1)r+3m}.
$$
The second word is obtained by exchanging the blocks, so its Gram matrix is
$$
 \mc G(p\,ab\,p\,ba\,p)=\mc G_p\,K^{\mathsf T}\mc G_p\,K\,\mc G_p ,
$$
and in the same way
$$
 K^{\mathsf T}\mc G_pK=\begin{pmatrix}9m+6r+h&-3m-r\\-3m-r&m\end{pmatrix},
 \qquad
 \mc G_pK^{\mathsf T}\mc G_pK=\begin{pmatrix}9m^2+3mr+1&-3m^2\\9mr+3r^2-3&1-3mr\end{pmatrix},
$$
and the first column of its Gram matrix is
$$
 \binom{(9m^2+3mr+1)m-3m^2r}{(9mr+3r^2-3)m+(1-3mr)r}
 =\binom{m(9m^2+1)}{(9m^2+1)r-3m}.
$$
The mean of the two roots is $(9m^2+1)r=Mr/m$. For $p=\operatorname{Pal}(d)$, \eqref{eq:literal-standard-identities}
gives $\phi_d(ba)=ba\,p$ and $\phi_d(ab)=ab\,p$, so that
$F_d(baab)=p\,\phi_d(ba)\phi_d(ab)=p\,ba\,p\,ab\,p$, and likewise for $abba$.
\end{proof}

Every row of \cref{tab:pairs} is \eqref{eq:one-block-closed-form}
evaluated at the ancestor of that row. The entries were also obtained
directly from the literal words, and the determinant, midpoint and count
identities are confirmed by exact integer arithmetic in the
supplementary material; the theorem, including its exhaustion statement,
depends on neither.

The theorem uses neither an odd-label hypothesis nor a factorization
of the common label. It also includes all three possible middle types
of a palindrome: empty, $a$, and $b$. Their apparent difference
disappears because the full-word decoder forces every reflected
preimage in the midpoint case to have the even form
\eqref{eq:classified-palindromes}. This is the main structural gain
of passing from square-root factors to whole Gram matrices.

\subsection{Every palindrome is a centre}
\label{subsec:any-centre}

\Cref{prop:one-block-closed-form} uses nothing about the centre beyond the symmetry and the
determinant of its Gram matrix, and the same is true of the construction half of
\cref{thm:palindromic-classification}. Fix a palindrome $p$ and put
$$
 G=\mu(p),\quad M_p=e^{\mathsf T}Ge,\quad \rho_p=(Ge)_1,\quad
 K_p=2ee^{\mathsf T}G-M_pI,\quad R_p=K_p/M_p .
$$
For blocks $W_1,\dots,W_n\in\{ab,ba\}$ write
$$
 x_p=p\,W_1\,p\,W_2\,p\cdots p\,W_n\,p ,
$$
and let $\widehat x_p$ be the word obtained by exchanging every block. The word $x_p$ is a
palindrome exactly when $n=2k$ is even and $W_{2k+1-i}$ is the exchange of $W_i$ for every
$i$, since a middle block would have to equal its own exchange. We then say that $x_p$ is
\emph{structured around} $p$, and $\widehat x_p$ is structured around $p$ as well.

\begin{proposition}[Block exchange around an arbitrary palindrome]
\label{prop:any-centre}
For every palindrome $p$ and all blocks $W_1,\dots,W_n\in\{ab,ba\}$,
$$
 \begin{gathered}
 K_p\,\mu(W_1p\cdots W_np)=\mu(\widehat W_1p\cdots\widehat W_np)\,K_p ,\\
 R_p\,\mu(W_1p\cdots W_np)\,e=\mu(\widehat W_1p\cdots\widehat W_np)\,e .
 \end{gathered}
$$
Moreover $e^{\mathsf T}\mu(x_p)e=e^{\mathsf T}\mu(\widehat x_p)e$; for palindromic block words
these are the Gram labels. If $x_p$ is structured around $p$ with $2k$ blocks and common label $M$, then
$\mu(\widehat x_p)=R_p^{\mathsf T}\mu(x_p)R_p$, the roots satisfy
$\rho(x_p)+\rho(\widehat x_p)=2M\rho_p/M_p$, and both words have completed counts
$(2k+1)\,c(p)$. For $k\geqslant1$ the two words are distinct and neither is central.
\end{proposition}

\begin{proof}
The derivation of \eqref{eq:decoder-intertwiners} uses only that $G$ is symmetric, through
$\tr(ABG)=3e^{\mathsf T}Ge$, and that the products have determinant one, through the
Cayley--Hamilton identity. With $G=\mu(p)$ it gives $K_p\,ABG=BAG\,K_p$ and
$K_p\,BAG=ABG\,K_p$. Hence $K_p\Pi=\widehat\Pi K_p$ for $\Pi=\mu(W_1p)\cdots\mu(W_np)$ and its
exchange $\widehat\Pi$, which is the displayed identity; applying it to $e$ and using
$K_pe=M_pe$ gives the column identity. Moreover $K_pe=2e(e^{\mathsf T}Ge)-M_pe=M_pe$ and
$e^{\mathsf T}GK_p=M_pe^{\mathsf T}G$. Since $\mu(x_p)=G\Pi$,
$$
 M_p\,e^{\mathsf T}\mu(\widehat x_p)e=e^{\mathsf T}G\widehat\Pi K_pe=e^{\mathsf T}GK_p\Pi e
 =M_p\,e^{\mathsf T}\mu(x_p)e .
$$
Next, $(ee^{\mathsf T}G)^2=M_pee^{\mathsf T}G$ gives $K_p^2=M_p^2I$, so
$\widehat\Pi=R_p\Pi R_p$; and $GR_pG^{-1}=R_p^{\mathsf T}$ because $G$ is symmetric. Therefore
$\mu(\widehat x_p)=GR_p\Pi R_p=R_p^{\mathsf T}\mu(x_p)R_p$. Since $P$ is symmetric, this gives
$\mc G(\widehat x_p)=(P^{-1}R_pP)^{\mathsf T}\,\mc G(x_p)\,(P^{-1}R_pP)$. Since $P^{-1}e=e_1$ and
$e^{\mathsf T}GP=e_1^{\mathsf T}PGP=(M_p,\rho_p)$,
$$
 P^{-1}R_pP=\frac{2e_1(M_p,\rho_p)-M_pI}{M_p}=\begin{pmatrix}1&2\rho_p/M_p\\0&-1\end{pmatrix},
$$
and the computation in the proof of \cref{lem:gram-bridge} shows that the congruence by this
matrix takes a Gram matrix with first column $(M,\rho)$ to one with first column
$(M,\,2M\rho_p/M_p-\rho)$. This is the root relation. The word $x_p$ contains $2k+1$ copies of
$p$ and $2k$ blocks, each with one letter of each kind, which gives the completed counts.
For $k\geqslant1$ their greatest common divisor is at least $2k+1$, while central words have
coprime completed counts by \cref{lem:standard-identities}; and the two words differ at the
first letter of the first block.
\end{proof}

The exchange also has a transparent arithmetic. Recall that the roots $r$ of $-1$ modulo a
label $M$ correspond bijectively to the primitive Gaussian integers $x+iy$ of norm $M$ taken
up to units, through $r\equiv xy^{-1}\pmod M$, and that replacing $r$ by $-r$ replaces $x+iy$
by its conjugate.

\begin{proposition}[The arithmetic of an exchange]
\label{prop:exchange-arithmetic}
Let $x_p$ be structured around $p$ with $2k\geqslant2$ blocks and common label $M$. Then
$M_p$ divides $M$, the quotient $N=M/M_p$ is prime to $M_p$, and
$$
 \rho(x_p)\equiv\rho(\widehat x_p)\equiv\rho_p\pmod{M_p},
 \qquad
 \rho(x_p)+\rho(\widehat x_p)\equiv0\pmod N .
$$
The exchange therefore keeps the Gaussian factor of norm $M_p$ and conjugates the factor of
norm $N$.
\end{proposition}

\begin{proof}
As in the proof of \cref{prop:one-block-closed-form}, inserting $P^{-1}P$ between the factors
gives $\mc G(x_p)=\mc G_pY_1\mc G_pY_2\cdots Y_n\mc G_p$, where $\mc G_p=\mc G(p)$ and $Y_i$ is
$K$ when $W_i=ba$ and $K^{\mathsf T}$ when $W_i=ab$. Directly from these two matrices,
$e_2^{\mathsf T}K=e_1^{\mathsf T}$, $e_2^{\mathsf T}K^{\mathsf T}=-e_1^{\mathsf T}$,
$Ke_2=-e_1$ and $K^{\mathsf T}e_2=e_1$; and modulo $M_p$, entrywise,
$e_1^{\mathsf T}\mc G_p\equiv\rho_pe_2^{\mathsf T}$ and $\mc G_pe_1\equiv\rho_pe_2$, because the
first row of $\mc G_p$ is $(M_p,\rho_p)$. Reading $M(x_p)=e_1^{\mathsf T}\mc G(x_p)e_1$ from the
left, each factor $\mc G_pY_i$ turns a multiple of $e_1^{\mathsf T}$ into $\pm\rho_p$ times
$e_1^{\mathsf T}$, so $M(x_p)\equiv\pm\rho_p^{\,n}\,e_1^{\mathsf T}\mc G_pe_1=\pm\rho_p^{\,n}M_p\equiv0$.
Reading $\rho(x_p)=e_2^{\mathsf T}\mc G(x_p)e_1$ from the right, the column starts as
$\rho_pe_2$, and each factor $\mc G_pY_i$ multiplies a multiple of $e_2$ by $\rho_p$ and by
$-1$ if $W_i=ba$, $+1$ if $W_i=ab$. Hence
$\rho(x_p)\equiv\rho_p^{\,2k+1}(-1)^{j}$, where $j$ is the number of blocks $ba$. For a
structured word the blocks pair off as $W_i$ and its exchange, so $j=k$; and
$\rho_p^2\equiv-1$ because $\det\mc G_p=1$. Hence $\rho(x_p)\equiv(-1)^k(-1)^k\rho_p=\rho_p$,
and the same holds for $\widehat x_p$.

For the coprimality write $\rho(x_p)=\rho_p+M_ps$ and $\rho(\widehat x_p)=\rho_p+M_ps'$. By
\cref{prop:any-centre} their sum is $2\rho_pN$, so $M_p(s+s')=2\rho_p(N-1)$. A prime dividing
both $M_p$ and $N$ divides neither $N-1$ nor $\rho_p$, which is prime to $M_p$, so it must
divide $2$; then $4$ divides $M=M_pN$. But $M$ is a sum of two coprime squares, the squared
norm of the first column of $T$, and such a sum is $1$ or $2$ modulo $4$. Hence
$\gcd(M_p,N)=1$, and $\rho(x_p)+\rho(\widehat x_p)=2\rho_pN\equiv0\pmod N$. By the Chinese
remainder theorem the two roots agree modulo $M_p$ and are opposite modulo $N$; a root and
its negative select conjugate Gaussian primes, which gives the last assertion.
\end{proof}

With the centre $p=baab$, whose label is $1130$ and root $467$, the closed form
\eqref{eq:one-block-closed-form} gives the label $1130\cdot11\,492\,101=12\,986\,074\,130$ and
the roots $5\,366\,814\,557$ and $5\,366\,807\,777$. Since $baab$ is not central, this pair lies
outside \cref{tab:pairs}, and the same label shows that the midpoint condition of
\cref{thm:palindromic-classification} cannot be removed.

\begin{remark}[The midpoint hypothesis cannot be dropped]
\label{rem:quadruple}
Exactly four palindromes carry the label of the example above,
$M=12\,986\,074\,130$, namely
$$
 \begin{aligned}
 Z_1&=baab\,ba\,baab\,ab\,baab, &\qquad Z_2&=baab\,ab\,baab\,ba\,baab,\\
 Z_3&=abba\,ba\,abba\,ab\,abba, &\qquad Z_4&=abba\,ab\,abba\,ba\,abba.
 \end{aligned}
$$
The pairs $(Z_1,Z_2)$ and $(Z_3,Z_4)$ are exchanges around $baab$ and $abba$. Read in blocks of
two letters, all four words lie in $\{ab,ba\}^*$, and $(Z_1,Z_4)$ and $(Z_2,Z_3)$ are exchanges
around the empty word, with eight blocks each. The remaining pairs $(Z_1,Z_3)$ and $(Z_2,Z_4)$
have equal labels and lie in the image of the empty directive, but neither word is the block
exchange of the other: their root midpoints are $2/5\pm3390/M$ instead of the ancestor
ratio $2/5$. Without the midpoint hypothesis, the conclusion of
\cref{thm:palindromic-classification} fails. A label can also be carried by more than two
palindromes, and two palindromes with the same label need not be exchanges of each other.
\end{remark}

The four words form a single class under exchanges, each around its own centre, and the two
pairs that are not exchanges have their midpoints shifted off every centre. The next subsection
shows that a midpoint at the ratio of a palindrome forces an exchange around it, at every centre,
and \cref{subsec:nested} explains the shifted midpoints.

\subsection{Decoding at every palindromic centre}
\label{subsec:centre-decoder}

\Cref{prop:any-centre} produces pairs of word columns related by the reflection $R_p$ of any
palindrome $p$. We now prove that there are no others, for every palindrome and all finite words,
and deduce the classification of \cref{thm:palindromic-classification} at every centre, without
the hypotheses of a common standard image and of a common prefix. The argument is shorter than that
of \cref{sec:decoder}. It works with the letters $a$ and $b$ in the slope coordinate, reads the
first two letters of each word, and then reads the whole centre letter by letter through the
transposed reflection. Beyond the symmetry of $\mu(p)$, the only number attached to the centre
that enters the estimates is the window in which its ratio lies; the letters of the centre are read
one at a time, by estimates valid for every word.

We write $t(v)=v_1/v_2$ for the slope of a column $v$, with $t(v)=\infty$ when $v_2=0$, and $t(w)$
for the slope of $\mu(w)e$; a matrix acts on slopes by fractional linear transformation. Two
nonzero columns span the same line exactly when their slopes agree. By \cref{sec:decoder}, every
word slope lies in $\mathcal J=[8/5,29/12]$, with $t(\varepsilon)=2$, $t(aw)\in A(\mathcal J)=[21/13,70/41]$
and $t(bw)\in B(\mathcal J)=[50/21,169/70]$, because $t(cw)=\mu(c)(t(w))$ and the two maps are
increasing on $\mathcal J$. So $t(w)<2$, $t(w)=2$ or $t(w)>2$ according as $w$ begins with $a$, is
empty, or begins with $b$, and no word slope is infinite, nonpositive, smaller than $21/13$ or
larger than $169/70$. We call this the \emph{first-letter rule}.

Fix a palindrome $p$ and write $G=\mu(p)$, $M=M_p$, $\rho=\rho_p$, $K=K_p$ and $\tau=2\rho/M$.
Since $Ge=(\rho,M-2\rho)^{\mathsf T}$ and $e^{\mathsf T}G=(Ge)^{\mathsf T}$,
\begin{equation}\label{eq:centre-K}
 K=\begin{pmatrix}4\rho-M&4M-8\rho\\2\rho&M-4\rho\end{pmatrix},\qquad
 K^{\mathsf T}=\begin{pmatrix}4\rho-M&2\rho\\4M-8\rho&M-4\rho\end{pmatrix}.
\end{equation}
Five facts follow by direct computation. First, $K^2=M^2I$, $Ke=Me$ and $GK=K^{\mathsf T}G$, as in
the proof of \cref{prop:any-centre}; hence $K^{\mathsf T}Ge=GKe=M\,Ge$, and
$K^{\mathsf T}(1,-2)^{\mathsf T}=2Ge\,e^{\mathsf T}(1,-2)^{\mathsf T}-M(1,-2)^{\mathsf T}=-M(1,-2)^{\mathsf T}$.
Second, multiplying \eqref{eq:centre-K} by $AB$, $BA$ and $A$,
\begin{equation}\label{eq:centre-intertwining}
 KAB=BAK^{\mathsf T},\qquad KBA=ABK^{\mathsf T},\qquad
 B^{-1}KA=M\begin{pmatrix}0&1\\1&-1-\tau\end{pmatrix};
\end{equation}
for instance both upper-left entries in the first identity equal $12(4\rho-M)+7(4M-8\rho)=16M-8\rho$.
Third, the slope $t(p)=\rho/(M-2\rho)$ lies in $\mathcal J$ and $\tau=2t(p)/(2t(p)+1)$ increases
with it, so $16/21\leqslant\tau\leqslant29/35$, the values at $8/5$ and at $29/12$. Fourth, for a
real $t$, with $\eta=1+\tau(t-2)$,
\begin{equation}\label{eq:centre-slopes}
 K\binom t1=M\binom{2\eta+2-t}{\eta},\qquad
 K^{\mathsf T}\binom t1=M\binom{(2\tau-1)t+\tau}{(4-4\tau)t+1-2\tau}.
\end{equation}
Fifth, a relation $[T\mu(x)e]=[\mu(y)e]$ with $T^2$ a nonzero multiple of $I$ is symmetric in $x$
and $y$: applying $T$ to both sides gives $[\mu(x)e]=[T\mu(y)e]$.

\begin{lemma}[The first two letters]
\label{lem:centre-first-letters}
Let $p$ be a palindrome and $x,y$ words with $[K_p\mu(x)e]=[\mu(y)e]$. Then $x$ is empty if and
only if $y$ is empty; if they are not empty, their first letters differ; and if $x$ begins with
$a$, then $x=ab\,x_1$ and $y=ba\,y_1$ for words $x_1,y_1$ with $[K_p^{\mathsf T}\mu(x_1)e]=[\mu(y_1)e]$.
\end{lemma}

\begin{proof}
\emph{First letters.} Let $t=t(x)$. Since $t-2\geqslant-2/5$ and $\tau\leqslant29/35$, the number
$\eta=1+\tau(t-2)$ is at least $1-58/175=117/175>0$. By \eqref{eq:centre-slopes} the column
$K\mu(x)e$ is a positive multiple of $(2\eta+2-t,\eta)^{\mathsf T}$, so
$t(y)=2+(2-t(x))/\eta$, and $t(y)-2$ has the sign of $2-t(x)$, zero included. The first-letter
rule gives the first two assertions.

\emph{The second letter of $x$.} Let $x=ax'$, so that $y=by'$. Multiplying
$[KA\mu(x')e]=[B\mu(y')e]$ by $B^{-1}$ and using \eqref{eq:centre-intertwining}, the column
$B^{-1}KA\mu(x')e$ is a positive multiple of $(1,\,t(x')-1-\tau)^{\mathsf T}$, so
$t(y')=1/(t(x')-1-\tau)$, read as $\infty$ when the denominator vanishes. If $x'$ were empty, then
$t(y')=1/(1-\tau)\geqslant21/5>169/70$. If $x'$ began with $a$, then
$t(x')\leqslant70/41<37/21\leqslant1+\tau$, and $t(y')$ would be negative. Neither is a word slope,
so $x'=bx_1$.

\emph{The second letter of $y$.} Now $[KAB\mu(x_1)e]=[B\mu(y')e]$, and
\eqref{eq:centre-intertwining} turns it into $[AK^{\mathsf T}\mu(x_1)e]=[\mu(y')e]$. By
\eqref{eq:centre-slopes}, with $t=t(x_1)\geqslant8/5$, the column $K^{\mathsf T}\mu(x_1)e$ is a
positive multiple of $((2\tau-1)t+\tau,\,(4-4\tau)t+1-2\tau)^{\mathsf T}$. Its first entry is
positive because $\tau>1/2$, and its second entry increases with $t$ and is at least
$(37-42\tau)/5\geqslant11/25$ at $t=8/5$. So $K^{\mathsf T}\mu(x_1)e$ has a positive slope $s$, and
$t(y')=A(s)$ lies strictly between $1$ and $2$. By the first-letter rule $y'=ay_1$, and cancelling
$A$ gives $[K^{\mathsf T}\mu(x_1)e]=[\mu(y_1)e]$.
\end{proof}

The transposed reflection $K^{\mathsf T}$ fixes the column $Ge$ of the centre and reverses
$(1,-2)^{\mathsf T}$. The next lemma shows that such a reflection can relate two word columns only if
both words begin as the centre does.

\begin{lemma}[One letter of a fixed word]
\label{lem:centre-one-letter}
Let $T$ be a real matrix and $\kappa\neq0$ with $Tu=\kappa u$ and $Tv=-\kappa v$, where $u=\mu(q)e$
for a nonempty word $q$ with first letter $c$, and $v$ is a column whose slope $\omega$ lies in
$[-1,-2/5]$. If $x$ and $y$ are words with $[T\mu(x)e]=[\mu(y)e]$, then both begin with $c$.
\end{lemma}

\begin{proof}
Put $\theta=t(q)$, which lies in $A(\mathcal J)$ if $c=a$ and in $B(\mathcal J)$ if $c=b$, and
$\Omega=-\omega\in[2/5,1]$. For a real $t\neq\omega$ let $\sigma(t)=(t-\theta)/(t-\omega)$. Writing
$(t,1)^{\mathsf T}=\alpha(\theta,1)^{\mathsf T}+\beta(\omega,1)^{\mathsf T}$ with
$\alpha=(t-\omega)/(\theta-\omega)$ and $\beta=(\theta-t)/(\theta-\omega)$, one has
$T(t,1)^{\mathsf T}=\kappa(\alpha(\theta,1)^{\mathsf T}-\beta(\omega,1)^{\mathsf T})$, and when
its slope $t'$ is finite, $\sigma(t')=\beta/\alpha=-\sigma(t)$. The relation gives $t(y)=t'$ for
$t=t(x)$, so $\sigma(t(y))=-\sigma(t(x))$. On $(\omega,\infty)$, which contains every word slope,
$\sigma$ is increasing, with derivative $(\theta-\omega)/(t-\omega)^2$, and $\sigma<1$.

\emph{The case $c=a$.} Suppose that $x$ does not begin with $a$, so $t(x)\geqslant2$ and
$\sigma(t(x))\geqslant\sigma_0=(2-\theta)/(2+\Omega)>0$. Solving $\sigma(t(y))=-\sigma(t(x))$ for
$t(y)$ and using that $\sigma\mapsto\sigma/(1+\sigma)$ increases,
$$
 t(y)=\theta-\frac{\sigma(t(x))}{1+\sigma(t(x))}(\theta+\Omega)
 \leqslant\theta-\frac{\sigma_0}{1+\sigma_0}(\theta+\Omega)
 =\theta-(2-\theta)\,\frac{\theta+\Omega}{4-\theta+\Omega}.
$$
The last quotient increases with $\Omega$, because $\theta<2$, and with $\theta$, so it is at least
its value $131/181$ at $\theta=21/13$, $\Omega=2/5$. With $\theta\leqslant70/41$ and
$2-\theta\geqslant12/41$ this gives
$t(y)\leqslant70/41-(12/41)(131/181)=11098/7421<21/13$, which is the slope of no word.

\emph{The case $c=b$.} Suppose that $x$ does not begin with $b$, so $t(x)\leqslant2<\theta$ and
$-\sigma(t(x))\geqslant\sigma_0=(\theta-2)/(2+\Omega)>0$; also $\sigma_0\leqslant29/168<1$. Since
$t(y)$ is a word slope, it lies in $(\omega,\infty)$, so $-\sigma(t(x))=\sigma(t(y))<1$. Solving as
before, and using that $\sigma\mapsto\sigma/(1-\sigma)$ increases on $[0,1)$,
$$
 t(y)\geqslant\theta+\frac{\sigma_0}{1-\sigma_0}(\theta+\Omega)
 =\theta+(\theta-2)\,\frac{\theta+\Omega}{4-\theta+\Omega}.
$$
The last quotient now decreases with $\Omega$, because $\theta>2$, and increases with $\theta$, so it
is at least its value $71/55$ at $\theta=50/21$, $\Omega=1$. With $\theta\geqslant50/21$ and
$\theta-2\geqslant8/21$ this gives $t(y)\geqslant50/21+(8/21)(71/55)=158/55>169/70$, which is the
slope of no word.

In both cases $x$ begins with $c$. The columns $u$ and $v$ have different slopes, so $T^2=\kappa^2I$,
the relation is symmetric, and the same argument shows that $y$ begins with $c$.
\end{proof}

\begin{lemma}[Reading the centre]
\label{lem:centre-reading}
Let $p$ be a palindrome. If two words $x$ and $y$ satisfy
$[K_p^{\mathsf T}\mu(x)e]=[\mu(y)e]$, then $x=p\,x_0$ and $y=p\,y_0$, where the words $x_0$ and
$y_0$ satisfy $[K_p\mu(x_0)e]=[\mu(y_0)e]$.
\end{lemma}

\begin{proof}
Write $p=c_1\cdots c_n$, and for $0\leqslant j\leqslant n$ put $\Pi_j=\mu(c_1\cdots c_j)$ and
$T_j=\Pi_j^{-1}K^{\mathsf T}\Pi_j$. Then $T_0=K^{\mathsf T}$ and $T_n=G^{-1}K^{\mathsf T}G=K$. For
$j<n$ let $q_j=c_{j+1}\cdots c_n$, a nonempty word with first letter $c_{j+1}$ and
$\Pi_j\mu(q_j)=G$. Then $T_j\mu(q_j)e=\Pi_j^{-1}K^{\mathsf T}Ge=M\mu(q_j)e$, and the column
$v_j=\Pi_j^{-1}(1,-2)^{\mathsf T}$ satisfies $T_jv_j=-Mv_j$. Its slope $\omega_j$ lies in
$[-1,-2/5]$. Indeed $\omega_0=-1/2$ and $\omega_j=\mu(c_j)^{-1}(\omega_{j-1})$; for $t<0$,
$$
 A^{-1}(t)=\frac{t-1}{2-t}\in\Bigl(-1,-\frac12\Bigr),\qquad
 B^{-1}(t)=\frac{t-2}{5-2t}\in\Bigl(-\frac12,-\frac25\Bigr),
$$
because $A^{-1}(t)+1=1/(2-t)>0$, $A^{-1}(t)+1/2=t/(2(2-t))<0$, $B^{-1}(t)+1/2=1/(2(5-2t))>0$ and
$B^{-1}(t)+2/5=t/(5(5-2t))<0$; induction on $j$ concludes.

Put $x^{(0)}=x$ and $y^{(0)}=y$, and suppose that $[T_j\mu(x^{(j)})e]=[\mu(y^{(j)})e]$ for some
$j<n$. \Cref{lem:centre-one-letter}, applied to $T_j$ with $q=q_j$ and $\kappa=M$, shows that
$x^{(j)}$ and $y^{(j)}$ both begin with $c_{j+1}$. Writing $x^{(j)}=c_{j+1}x^{(j+1)}$ and
$y^{(j)}=c_{j+1}y^{(j+1)}$ and multiplying by $\mu(c_{j+1})^{-1}$ gives
$[T_{j+1}\mu(x^{(j+1)})e]=[\mu(y^{(j+1)})e]$. After $n$ steps $x=p\,x^{(n)}$ and $y=p\,y^{(n)}$, and
the relation for $T_n=K$ is the claim with $x_0=x^{(n)}$ and $y_0=y^{(n)}$.
\end{proof}

\begin{theorem}[Decoding at every palindromic centre]
\label{thm:centre-decoder}
For every palindrome $p$ and all finite words $x$ and $y$, the empty word included, the lines
$[R_p\mu(x)e]$ and $[\mu(y)e]$ coincide if and only if $x\in\{ab\,p,\ ba\,p\}^*$ and $y$ is
obtained from $x$ by exchanging every block. In that case $R_p\mu(x)e=\mu(y)e$ exactly.
\end{theorem}

\begin{proof}
The implication from right to left and the exact equality are \cref{prop:any-centre}: for a block
word $x$ and its exchange $\widehat x$, $K\mu(x)=\mu(\widehat x)K$, and $Ke=Me$. For the converse we
argue by induction on $\abs x+\abs y$. Suppose $[K\mu(x)e]=[\mu(y)e]$. If one of the words is empty,
both are, by \cref{lem:centre-first-letters}, and the empty word is a block word equal to its
exchange. Otherwise their first letters differ. If $x$ begins with $a$,
\cref{lem:centre-first-letters} gives $x=ab\,x_1$, $y=ba\,y_1$ and
$[K^{\mathsf T}\mu(x_1)e]=[\mu(y_1)e]$, and \cref{lem:centre-reading} gives $x_1=p\,x''$,
$y_1=p\,y''$ and $[K\mu(x'')e]=[\mu(y'')e]$. Since $\abs{x''}+\abs{y''}<\abs x+\abs y$, the
induction hypothesis shows that $x''$ is a block word and $y''$ its exchange, so $x=ab\,p\,x''$ is a
block word and $y=ba\,p\,y''$ is its exchange. If $x$ begins with $b$, the relation is symmetric,
because $K^2=M^2I$; the previous case applied to $(y,x)$ gives $y=ab\,p\,y''$ and $x=ba\,p\,x''$ with
$[K\mu(y'')e]=[\mu(x'')e]$, hence $[K\mu(x'')e]=[\mu(y'')e]$, and the induction hypothesis concludes
in the same way.
\end{proof}

For a standard centre $p=\operatorname{Pal}(d)$ one has $ab\,p=\phi_d(ab)$ and $ba\,p=\phi_d(ba)$ by
\eqref{eq:literal-standard-identities}, so the theorem contains \cref{thm:word-decoder} and shows
moreover that no word outside the image of $\phi_d$ is related to a word column by $R_d$. The proof
uses neither the standard letters $U$, $V$ nor the cylinders of \cref{sec:decoder}; those cylinders
are still needed in \cref{sec:displaced}, where they measure how far a column moves when the
midpoint is displaced. The theorem has the following consequence for palindromes.

\begin{corollary}[Classification at every centre]
\label{cor:centre-classification}
Let $p$, $Z$ and $Z'$ be palindromes. Then $M(Z)=M(Z')$ and $\rho(Z)+\rho(Z')=2M(Z)\rho_p/M_p$ if and
only if $Z=p\,X$ and $Z'=p\,\widehat X$ for a block word $X\in\{ab\,p,\ ba\,p\}^*$. In that case either
$Z=Z'=p$, or $Z$ is structured around $p$ and $Z'$ is its exchange, and the two words have the same
numbers of letters $a$ and $b$. In particular two distinct central words never have equal labels
with root midpoint at the ratio of a palindrome.
\end{corollary}

\begin{proof}
Put $J_\tau=\left(\begin{smallmatrix}1&\tau\\0&-1\end{smallmatrix}\right)$ with $\tau=2\rho_p/M_p$;
by the proof of \cref{prop:any-centre}, $P^{-1}R_pP=J_\tau$. The computation in the proof of
\cref{lem:gram-bridge} shows that $J_\tau^{\mathsf T}\mathcal G(Z)J_\tau$ is symmetric of
determinant one with first row $(M(Z),\,\tau M(Z)-\rho(Z))$, and a symmetric matrix of determinant
one with positive upper-left entry is determined by its first row. So the two conditions on labels
and roots hold if and only if $\mathcal G(Z')=J_\tau^{\mathsf T}\mathcal G(Z)J_\tau$, that is, since $P$
is symmetric, $\mu(Z')=R_p^{\mathsf T}\mu(Z)R_p$. If this holds, then
$\mu(Z')e=R_p^{\mathsf T}\mu(Z)e$ because $R_pe=e$, and \cref{lem:centre-reading} gives $Z=p\,X$ and
$Z'=p\,Y$ with $[K_p\mu(X)e]=[\mu(Y)e]$; by \cref{thm:centre-decoder}, $X$ is a block word and
$Y=\widehat X$. Conversely, if $Z=p\,X$ and $Z'=p\,\widehat X$, then $K_p\mu(X)=\mu(\widehat X)K_p$ and
$K_p^2=M_p^2I$ give $\mu(\widehat X)=R_p\mu(X)R_p$, so
$\mu(Z')=G\,R_p\mu(X)R_p=R_p^{\mathsf T}G\mu(X)R_p=R_p^{\mathsf T}\mu(Z)R_p$. Since $Z$ is a palindrome
and every block has length $\abs p+2$, comparing $Z$ with its reversal shows that $X$ has an even
number $2k$ of blocks with $W_{2k+1-i}=\widehat W_i$; so either $k=0$, or $Z$ is structured around
$p$ and $Z'$ is its exchange. Exchanging blocks does not change the letter counts. Finally a word
structured with $k\geqslant1$ has completed counts $(2k+1)\,c(p)$, which are not coprime, so it is
not central.
\end{proof}

The centre in the corollary is determined by the pair: the fraction $\rho_p/M_p$ is reduced,
because $\rho_p^2\equiv-1\pmod{M_p}$, so it gives $M_p$ and $\rho_p$, hence the Gram matrix of $p$
and, the Cohn coding being injective, the palindrome $p$ itself. The corollary contains
\cref{thm:palindromic-classification}: when $Z=F_d(z_X)$ and $Z'=F_d(z_Y)$ and $p=\operatorname{Pal}(d)$,
the block word $X$ lies in $\phi_d(\{ab,ba\}^*)$, and the injectivity of $F_d$ returns the pair
$(\widetilde ww,\widetilde{\widehat w}\,\widehat w)$. \Cref{rem:quadruple} shows that the hypothesis on
the midpoint cannot be dropped.

\subsection{Nested centres and shifted midpoints}
\label{subsec:nested}

The quadruple of \cref{rem:quadruple} raises two questions: why do four palindromes share a
label, and why are the midpoints of two of its pairs shifted off every centre? The answer is that
exchanges around nested centres compose. When the centre of an exchange is itself one of two
exchanged palindromes, the words built around the two of them form a square of twins, and nesting
$r$ levels gives a cube with $2^r$ vertices. The shift of a midpoint is then the geometry of two
reflections of the ratio axis, which compose to a translation.

It is convenient to name the words of \cref{subsec:any-centre} by their blocks. A \emph{block
sequence} is a finite sequence $B=(B_1,\dots,B_n)$ with $B_i\in\{ab,ba\}$; we write
$\widehat B=(\widehat B_1,\dots,\widehat B_n)$, and we call $B$ \emph{structured} when $n=2k\geqslant2$
and $B_{n+1-i}=\widehat B_i$ for every $i$. For a word $q$ put
$$
 q\langle B\rangle=q\,B_1\,q\,B_2\,q\cdots q\,B_n\,q .
$$
The words structured around a palindrome $p$ are exactly the words $p\langle B\rangle$ with $B$
structured, and the exchange of $p\langle B\rangle$ around $p$ is $p\langle\widehat B\rangle$. For
block sequences $X=(X_1,\dots,X_n)$ and $Y=(Y_1,\dots,Y_m)$ let
$$
 X\ast Y=(X,\,Y_1,\,X,\,Y_2,\,\dots,\,X,\,Y_m,\,X)
$$
be the block sequence in which the blocks of $X$ are written $m+1$ times, separated by
$Y_1,\dots,Y_m$; it has $(m+1)n+m$ blocks.

\begin{lemma}[Reading through an inner centre]
\label{lem:nested-reading}
For every word $q$ and all block sequences $X$ and $Y$,
$$
 \bigl(q\langle X\rangle\bigr)\langle Y\rangle=q\langle X\ast Y\rangle,
 \qquad \widehat{X\ast Y}=\widehat X\ast\widehat Y .
$$
If $X$ and $Y$ are structured, then so is $X\ast Y$.
\end{lemma}

\begin{proof}
The word $Q=q\langle X\rangle$ begins and ends with $q$. In $Q\langle Y\rangle=Q\,Y_1\,Q\cdots Y_m\,Q$
every junction reads $\cdots X_n\,q\,Y_j\,q\,X_1\cdots$, so the whole word is $q$ followed by the
blocks $X_1,\dots,X_n,Y_1,X_1,\dots,X_n,Y_2,\dots,Y_m,X_1,\dots,X_n$, each of them followed by one
copy of $q$. This is $q\langle X\ast Y\rangle$. The second identity holds block by block. For the
last assertion let $X$ have $n=2k$ blocks and $Y$ have $m=2j$ blocks. Reading $X\ast Y$ backwards
gives $(\overleftarrow X,Y_m,\overleftarrow X,\dots,Y_1,\overleftarrow X)$, where
$\overleftarrow X=(X_n,\dots,X_1)$. Structure of $X$ means $\overleftarrow X=\widehat X$, and
structure of $Y$ means $Y_{m+1-i}=\widehat Y_i$, so the reversed sequence is
$(\widehat X,\widehat Y_1,\widehat X,\dots,\widehat Y_m,\widehat X)=\widehat X\ast\widehat Y$, which
is $\widehat{X\ast Y}$ by the second identity. Its number of blocks, $(m+1)n+m$, is even and at
least two.
\end{proof}

\begin{theorem}[Cubes of nested exchanges]
\label{thm:nested-cubes}
Let $p$ be a palindrome, let $r\geqslant1$, and let $B_1,\dots,B_r$ be structured block sequences,
$B_i$ with $2k_i$ blocks. For $s\in\{0,1\}^r$ put $B_i^s=B_i$ if $s_i=0$ and $B_i^s=\widehat B_i$
if $s_i=1$, and define
$$
 Q_0^s=p,\qquad Q_i^s=Q_{i-1}^s\langle B_i^s\rangle\quad(1\leqslant i\leqslant r),\qquad Z^s=Q_r^s .
$$
For $1\leqslant i\leqslant r$ let $f_i\in\{0,1\}^r$ have its entries equal to one exactly in the
positions $i,\dots,r$, and write $s+f_i$ for the sum modulo two. Then:
\begin{enumerate}[label=\textup{(\roman*)}]
\item the $2^r$ words $Z^s$ are pairwise distinct palindromes, all with completed counts
 $(2k_1+1)\cdots(2k_r+1)\,c(p)$;
\item for every $s$ and every $i$, the words $Z^s$ and $Z^{s+f_i}$ are exchanges of each other
 around the palindrome $Q_{i-1}^s$;
\item all $2^r$ words have the same Gram label $M$, any two of them are joined by at most $r$
 exchanges of the kind in \textup{(ii)}, and the roots of the pair in \textup{(ii)} satisfy
 $\rho(Z^s)+\rho(Z^{s+f_i})=2M\rho(Q_{i-1}^s)/M(Q_{i-1}^s)$.
\end{enumerate}
\end{theorem}

\begin{proof}
\emph{Every level is a centre.} Define block sequences $R_r^s=B_r^s$ and
$R_i^s=B_i^s\ast R_{i+1}^s$ for $1\leqslant i<r$. We claim that $Z^s=Q_{i-1}^s\langle R_i^s\rangle$
for every $i$, and prove it by downward induction on $i$. For $i=r$ it is the definition
$Z^s=Q_r^s=Q_{r-1}^s\langle B_r^s\rangle$. If $Z^s=Q_i^s\langle R_{i+1}^s\rangle$, then
$Q_i^s=Q_{i-1}^s\langle B_i^s\rangle$ and \cref{lem:nested-reading} give
$Z^s=Q_{i-1}^s\langle B_i^s\ast R_{i+1}^s\rangle=Q_{i-1}^s\langle R_i^s\rangle$. By the last
assertion of \cref{lem:nested-reading} and induction, every $R_i^s$ is structured. The same
argument applied to the word $Q_{i-1}^s$ in place of $Z^s$ writes it as $p\langle R\rangle$ with $R$
structured, or as $p$ itself when $i=1$; in both cases $Q_{i-1}^s$ is a palindrome. Hence $Z^s$ is
structured around the palindrome $Q_{i-1}^s$.

\emph{The exchange is a flip of the last coordinates.} The exchange of $Z^s$ around $Q_{i-1}^s$ is
$Q_{i-1}^s\langle\widehat{R_i^s}\rangle$. Applying the second identity of \cref{lem:nested-reading}
repeatedly, $\widehat{R_i^s}=\widehat{B_i^s}\ast\widehat{R_{i+1}^s}=\cdots$ is the sequence $R_i^{s'}$
built from $s'=s+f_i$, the vector obtained by changing $s_i,\dots,s_r$. The first $i-1$ coordinates
of $s'$ are those of $s$, so $Q_{i-1}^{s'}=Q_{i-1}^s$, and the exchange is
$Q_{i-1}^{s'}\langle R_i^{s'}\rangle=Z^{s'}$. This proves (ii).

\emph{Palindromes, distinctness and counts.} The case $i=1$ of the first step gives
$Z^s=p\langle R_1^s\rangle$ with $R_1^s$ structured, so $Z^s$ is a palindrome. Let $s\neq s'$, let
$i$ be the first coordinate in which they differ, and put $Q=Q_{i-1}^s=Q_{i-1}^{s'}$. Every
$Q_j^s$ begins with $Q_{j-1}^s$, so $Z^s$ begins with $Q\langle B_i^s\rangle$ and $Z^{s'}$ begins
with $Q\langle\widehat{B_i^s}\rangle$; the two words differ in the letter that follows the first
copy of $Q$, which is the first letter of a block in one and of its exchange in the other. By
\cref{prop:any-centre}, a word structured around a centre with $2k$ blocks has completed counts
$2k+1$ times those of the centre; applied at the levels $1,\dots,r$ this gives the counts in (i).

\emph{Labels and roots.} By \cref{prop:any-centre}, two words exchanged around a palindrome have the
same label, and their roots sum to twice the label times the ratio of the centre. The vectors
$f_1,\dots,f_r$ are linearly independent modulo two, since the matrix with these rows is triangular
with ones on the diagonal. Every vector of $\{0,1\}^r$ is therefore a sum of at most $r$ of them,
and applying the corresponding exchanges one after the other joins $Z^s$ to any other $Z^{s'}$.
The label is therefore common to all $2^r$ words, which proves (iii).
\end{proof}

For $r=2$ the cube is a square, and its two pairs that are not exchanges are the pairs of the
quadruple whose midpoints are shifted.

\begin{corollary}[The square and its shifted midpoints]
\label{cor:nested-square}
Let $p$ be a palindrome and $W$, $V$ structured block sequences; put $P_1=p\langle W\rangle$,
$P_2=p\langle\widehat W\rangle$ and
$$
 Z_1=P_1\langle V\rangle,\qquad Z_2=P_1\langle\widehat V\rangle,\qquad
 Z_3=P_2\langle V\rangle,\qquad Z_4=P_2\langle\widehat V\rangle ,
$$
with common label $M$ and roots $\rho_1,\dots,\rho_4$. The pairs $(Z_1,Z_2)$ and $(Z_3,Z_4)$ are
exchanges around $P_1$ and $P_2$, and the pairs $(Z_1,Z_4)$ and $(Z_2,Z_3)$ are exchanges around
$p$. The two remaining pairs satisfy $\rho_1\neq\rho_2$ and
$$
 \frac{\rho_1+\rho_3}{2M}=\frac{\rho_p}{M_p}+\frac{\rho_1-\rho_2}{2M},\qquad
 \frac{\rho_2+\rho_4}{2M}=\frac{\rho_p}{M_p}-\frac{\rho_1-\rho_2}{2M}.
$$
If $W$ and $V$ have two blocks each, then $M(P_1)=M_p(9M_p^2+1)$, $M=M(P_1)\bigl(9M(P_1)^2+1\bigr)$
and $\abs{\rho_1-\rho_2}=6M(P_1)$.
\end{corollary}

\begin{proof}
Take $r=2$, $B_1=W$ and $B_2=V$ in \cref{thm:nested-cubes}; then $Z_1,Z_2,Z_3,Z_4$ are
$Z^{(0,0)},Z^{(0,1)},Z^{(1,0)},Z^{(1,1)}$. The flip $f_2=(0,1)$ joins $Z_1$ to $Z_2$ around
$Q_1^{(0,0)}=P_1$ and $Z_3$ to $Z_4$ around $Q_1^{(1,0)}=P_2$. The flip $f_1=(1,1)$ joins $Z_1$ to
$Z_4$ and $Z_2$ to $Z_3$, both around $Q_0=p$. By part (iii) of the theorem,
$\rho_1+\rho_4=2M\rho_p/M_p$ and $\rho_2+\rho_3=2M\rho_p/M_p$. Hence
$$
 \rho_1+\rho_3=\rho_1-\rho_2+2M\frac{\rho_p}{M_p},\qquad
 \rho_2+\rho_4=\rho_2-\rho_1+2M\frac{\rho_p}{M_p},
$$
which are the displayed identities after division by $2M$. The words $Z_1$ and $Z_2$ are distinct
palindromes with the same label, and a palindrome is determined by its label and root, because
they determine the Gram matrix through its determinant and the Cohn coding is injective; so
$\rho_1\neq\rho_2$. If $W$ and $V$ have two blocks each, \cref{prop:one-block-closed-form}
applied around $p$ gives the label of $P_1$, and applied around $P_1$ it gives $M$ and
$\rho_1-\rho_2=\pm6M(P_1)$.
\end{proof}

For $p=\varepsilon$ and $W=V=(ba,ab)$ one has $P_1=baab$, $P_2=abba$, $M(P_1)=1130$, and the four
words are those of \cref{rem:quadruple}, with $M=1130\cdot11\,492\,101$. The two pairs that are not
exchanges have their midpoints at $2/5\pm3390/M$: the shift is half the width $6\cdot1130$ of the
rungs around $baab$ and $abba$, divided by the label. \Cref{fig:nested-square} draws this square
above the rung of $abba$ and $baab$.

\begin{figure}[htbp]
\centering
% Author: Dr. Denys Dutykh (Khalifa University of Science and Technology, Abu Dhabi, UAE)
% Generated by supplement/code/make_figures.py from exact data; do not edit by hand.

\begin{tikzpicture}[font=\footnotesize]
\draw[ink, line width=0.5pt] (0,0) -- (10.8,0);
\draw[ink, line width=0.5pt] (1.35,0) -- (1.35,-0.08);
\node[ink, below] at (1.35,-0.08) {$0.385$};
\draw[ink, line width=0.5pt] (2.7,0) -- (2.7,-0.08);
\node[ink, below] at (2.7,-0.08) {$0.390$};
\draw[ink, line width=0.5pt] (4.05,0) -- (4.05,-0.08);
\node[ink, below] at (4.05,-0.08) {$0.395$};
\draw[ink, line width=0.5pt] (5.4,0) -- (5.4,-0.08);
\node[ink, below] at (5.4,-0.08) {$0.400$};
\draw[ink, line width=0.5pt] (6.75,0) -- (6.75,-0.08);
\node[ink, below] at (6.75,-0.08) {$0.405$};
\draw[ink, line width=0.5pt] (8.1,0) -- (8.1,-0.08);
\node[ink, below] at (8.1,-0.08) {$0.410$};
\draw[ink, line width=0.5pt] (9.45,0) -- (9.45,-0.08);
\node[ink, below] at (9.45,-0.08) {$0.415$};
\node[ink, below] at (5.4,-0.48) {$\rho/M$};
\node[ink, left] at (-0.1,0.55) {$M=5$};
\node[ink, left] at (-0.1,2.35) {$M=1130$};
\node[ink, left] at (-0.1,4.75) {$M=12\,986\,074\,130$};
\draw[slate, densely dotted, line width=0.4pt] (5.4,0.55) -- (5.4,4.75);
\draw[slate, densely dotted, line width=0.4pt] (8.984,2.35) -- (8.984,4.75);
\draw[slate, densely dotted, line width=0.4pt] (1.816,2.35) -- (1.816,4.75);
\draw[navy, line width=0.8pt] (1.816,2.35) -- (8.984,2.35);
\draw[navy, line width=0.8pt] (2.366,4.75) .. controls (2.366,5.5) and (8.434,5.5) .. (8.434,4.75);
\draw[navy, line width=0.8pt] (1.266,4.75) .. controls (1.266,5.9) and (9.534,5.9) .. (9.534,4.75);
\draw[wine, line width=0.9pt] (8.434,4.75) -- (9.534,4.75);
\draw[wine, line width=0.9pt] (1.266,4.75) -- (2.366,4.75);
\draw[ochre, fill=white, line width=0.6pt] (4.85,4.43) circle (1.5pt);
\draw[ochre, fill=white, line width=0.6pt] (5.95,4.43) circle (1.5pt);
\draw[slate, line width=0.35pt] (4.8,4.37) -- (4.436,4.13);
\node[ink, right, align=left, inner sep=1pt] at (2.136,3.43) {midpoints $2/5\pm3/N$\\of the diagonals\\$(Z_1,Z_3)$ and $(Z_2,Z_4)$};
\draw[wine, fill=white, line width=0.45pt] (9.534,4.75) circle (1.3pt);
\node[ink, below] at (9.534,4.69) {$Z_1$};
\draw[wine, fill=white, line width=0.45pt] (8.434,4.75) circle (1.3pt);
\node[ink, below] at (8.434,4.69) {$Z_2$};
\draw[wine, fill=white, line width=0.45pt] (2.366,4.75) circle (1.3pt);
\node[ink, below] at (2.366,4.69) {$Z_3$};
\draw[wine, fill=white, line width=0.45pt] (1.266,4.75) circle (1.3pt);
\node[ink, below] at (1.266,4.69) {$Z_4$};
\filldraw[ochre, draw=ink, line width=0.25pt] (5.4,0.55) +(0,1.9pt) -- +(1.9pt,0) -- +(0,-1.9pt) -- +(-1.9pt,0) -- cycle;
\node[ink, right=2.2pt] at (5.4,0.55) {$\varepsilon$};
\filldraw[ochre, draw=ink, line width=0.25pt] (8.984,2.35) +(0,1.9pt) -- +(1.9pt,0) -- +(0,-1.9pt) -- +(-1.9pt,0) -- cycle; \draw[wine, line width=0.45pt] (8.984,2.35) circle (2.7pt);
\node[ink, below right=1pt] at (8.984,2.35) {$baab$};
\filldraw[ochre, draw=ink, line width=0.25pt] (1.816,2.35) +(0,1.9pt) -- +(1.9pt,0) -- +(0,-1.9pt) -- +(-1.9pt,0) -- cycle; \draw[wine, line width=0.45pt] (1.816,2.35) circle (2.7pt);
\node[ink, below left=1pt] at (1.816,2.35) {$abba$};
\draw[navy, line width=0.8pt] (0.2,6.8) -- (0.8,6.8); \node[ink, right] at (0.85,6.8) {exchange around $\varepsilon$};
\draw[wine, line width=0.9pt] (4.6,6.8) -- (5.2,6.8); \node[ink, right] at (5.25,6.8) {exchange around $baab$ or $abba$};
\filldraw[ochre, draw=ink, line width=0.25pt] (0.5,6.35) +(0,1.9pt) -- +(1.9pt,0) -- +(0,-1.9pt) -- +(-1.9pt,0) -- cycle; \node[ink, right] at (0.85,6.35) {centre of an exchange};
\draw[ochre, fill=white, line width=0.6pt] (4.9,6.35) circle (1.5pt); \node[ink, right] at (5.25,6.35) {midpoint of a diagonal};
\end{tikzpicture}

\caption{The square of \cref{cor:nested-square} for $p=\varepsilon$ and $W=V=(ba,ab)$, in the plane
of the ratio $\rho/M$ and the label, at the labels $5$, $1130$ and $12\,986\,074\,130$. The two lower
levels are drawn to scale. At the top level $Z_1$ and $Z_2$ lie at $467/1130\pm3/N$ and $Z_3$ and
$Z_4$ at $437/1130\pm3/N$, where $N=9\cdot1130^2+1=11\,492\,101$; these offsets, about
$2.6\cdot10^{-7}$, are drawn enlarged, all by the same factor. The short rungs are the exchanges
around $baab$ and $abba$, each centred above its centre; the arcs are the exchanges around the empty
word, centred at $2/5$ like the rung below them. The diagonals $(Z_1,Z_3)$ and $(Z_2,Z_4)$ are not
exchanges: their midpoints $2/5\pm3/N$ are the inner centre shifted by half the width of a short
rung.}
\label{fig:nested-square}
\end{figure}

When every block sequence has two blocks, the arithmetic of a cube is completely explicit. Put
$M_i=M(Q_i^s)$, which by \cref{prop:one-block-closed-form} does not depend on $s$, and
$N_i=9M_i^2+1$ for $0\leqslant i<r$, so that $M_0=M_p$ and $M_i=M_{i-1}N_{i-1}$. The same
proposition gives $\rho(Q_i^s)=N_{i-1}\rho(Q_{i-1}^s)\pm3M_{i-1}$, with the sign $+$ exactly when the
first block of $B_i^s$ is $ba$. Since $N_{i-1}\equiv1$ modulo $M_{i-1}=M_pN_0\cdots N_{i-2}$, the
numbers $M_p,N_0,\dots,N_{r-1}$ are pairwise coprime with product $M$. Reducing the recursion,
$\rho(Q_i^s)\equiv\rho(Q_{i-1}^s)$ modulo $M_{i-1}$ and $\rho(Q_i^s)\equiv\pm3M_{i-1}$ modulo
$N_{i-1}$; since every later level keeps the residue modulo the label of the previous one,
$$
 \rho(Z^s)\equiv\rho_p\pmod{M_p},\qquad \rho(Z^s)\equiv\pm3M_{i-1}\pmod{N_{i-1}}\quad(1\leqslant i\leqslant r),
$$
and different $s$ give different patterns of signs. In the correspondence between roots of $-1$
and Gaussian integers recalled before \cref{prop:exchange-arithmetic}, $N_{i-1}=(3M_{i-1})^2+1$ is
the norm of $3M_{i-1}\pm i$, whose roots are $\pm3M_{i-1}$. The $2^r$ twins of a cube therefore carry
the Gaussian integers $\gamma_p\prod_{i=1}^r(3M_{i-1}\pm i)$, with $\gamma_p$ the Gaussian integer
attached to $\rho_p$, for all
choices of signs, and the exchange around $Q_{i-1}^s$ conjugates the factors of the levels
$i,\dots,r$. This is the sign change of \cref{prop:exchange-arithmetic}, made explicit. Whether every
coincidence of labels among palindromes arises from such cubes is the question of
\cref{sec:outlook}.
